% Generated draft: requires mathematical and source review.
\documentclass[11pt]{article}
\usepackage[margin=25mm]{geometry}
\usepackage{fontspec}
\setmainfont{Latin Modern Roman}
\usepackage{amsmath,amssymb,mathtools,unicode-math}
\setmathfont{Latin Modern Math}
\usepackage{enumitem,needspace,xurl,hyperref}
\setlength{\parindent}{0pt}
\setlength{\parskip}{5pt}
\setlength{\emergencystretch}{3em}
\newcommand{\N}{\mathbb{N}}
\newcommand{\Z}{\mathbb{Z}}
\newcommand{\Q}{\mathbb{Q}}
\newcommand{\R}{\mathbb{R}}
\newcommand{\C}{\mathbb{C}}
\newcommand{\F}{\mathbb{F}}
\newcommand{\A}{\mathbb{A}}
\newcommand{\PP}{\mathbb{P}}
\newcommand{\E}{\mathbb{E}}
\newcommand{\eps}{\varepsilon}
\newcommand{\dd}{\,\mathrm{d}}
\begin{document}
\section*{UnsolvedMath ID: 701359}
\textbf{Source problem code:} catalogue entry
\subsection{Definitions and mathematical statement}
For each positive integer \(m\) and integer \(r\) with \(0\le r<m\), define the residue class
\[
r(m)=\{r+km:k\in\Z\}.
\]
For disjoint residue classes \(r_1(m_1)\) and \(r_2(m_2)\), define their class transposition \(\tau_{r_1(m_1),r_2(m_2)}\in\operatorname{Sym}(\Z)\) by
\[
\tau_{r_1(m_1),r_2(m_2)}(n)=
\begin{cases}
r_2+km_2,&n=r_1+km_1,\
r_1+km_1,&n=r_2+km_2,\
n,&n\notin r_1(m_1)\cup r_2(m_2),
\end{cases}
\qquad(k\in\Z).
\]
Let
\[
G=\mathrm{CT}(\Z)=\left\langle
\tau_{r_1(m_1),r_2(m_2)}:
\begin{gathered}
m_1,m_2\in\Z_{>0},\quad 0\le r_i<m_i\ (i=1,2),\\
r_1(m_1)\cap r_2(m_2)=\varnothing
\end{gathered}
\right\rangle
\le\operatorname{Sym}(\Z).
\]
The supplied source states that \(G\) is simple and cites \emph{Math. Z.} \textbf{264}, no.~4 (2010), 927--938 for this background assertion.

Let \(\rho\in\operatorname{Sym}(\Z)\) be the reflection \(\rho(n)=-n-1\), and let \(\alpha\in\operatorname{Aut}(G)\) be the automorphism induced by conjugation,
\[
\alpha(\sigma)=\rho^{-1}\circ\sigma\circ\rho\qquad(\sigma\in G).
\]
Write \(\operatorname{Inn}(G)\) for the subgroup of inner automorphisms, set \(\operatorname{Out}(G)=\operatorname{Aut}(G)/\operatorname{Inn}(G)\), and denote the image of \(\alpha\) in this quotient by \([\alpha]\).

\textbf{Question (S. Kohl, 2010).} Is
\[
\operatorname{Out}(\mathrm{CT}(\Z))=\langle[\alpha]\rangle\cong\mathrm C_2,
\]
where \(\mathrm C_2\) is the cyclic group of order two? Equivalently, is \(\alpha\) non-inner, and does every \(\beta\in\operatorname{Aut}(G)\) admit a representation
\[
\beta=\iota_g\circ\alpha^{\varepsilon}
\quad\text{for some }g\in G\text{ and }\varepsilon\in\{0,1\},
\qquad \iota_g(x)=gxg^{-1}?
\]
The source's generator notation is interpreted as the outer class of the displayed conjugation automorphism. The supplied registry labels the problem open; this label is not an independently verified literature assessment.
\subsection{Short English statement}
S. Kohl asks whether every automorphism of the group generated by class transpositions of \(\Z\) is an inner automorphism composed with either the identity or conjugation by \(n\mapsto -n-1\), and whether this reflection induces a nontrivial outer automorphism. The supplied open status requires independent verification.
\subsection{Sources}
\begin{itemize}
\item Local UnsolvedMath registry: permanent ID \(701359\), exact source code \texttt{catalogue entry}, title \emph{Kourovka Notebook Problem 17.57}. The supplied record attributes the problem to S. Kohl (2010) and labels it open; its update date is \texttt{2026-10-16}.
\item \emph{The Kourovka Notebook (Unsolved Problems in Group Theory)}, original problem 17.57, 17th issue (2010); also identified in the supplied record with the 21st issue edition (2026), arXiv:1401.0300v48. The supplied collection description names E. I. Khukhro and V. D. Mazurov as editors.
\item Background simplicity reference, as supplied: \emph{Math. Z.} \textbf{264}, no.~4 (2010), 927--938. The record does not supply the article title or authors.
\end{itemize}
\end{document}
